%% =====================================================================
%% 第 21 章  幂次修正、次领头幂次、重整化子与非微扰地位
%% (WRITER_SPEC_v2 §1.1: key power, file ch21_power.tex; writer W24)
%% =====================================================================
\chapter{幂次修正、次领头幂次、重整化子与非微扰地位}\label{ch:power}

前面各章的因子化定理都是渐近展开的领头项加一个“被幂次压低”的余项。本章讨论余项本身：它的结构、它在什么意义下可以被控制，以及“精确到幂次修正”在 QCD 中究竟是什么意思。领头幂次之外的严格性远不如领头幂次整齐：DIS 矩的扭度 4 修正有以 OPE 为前提的严格算符结构（\cref{prop:dis-twist4}），其 $x$ 空间形式只有全阶论证；单自旋不对称的扭度 3 因子化只在有限阶得到验证；次领头幂次（NLP）的因子化因端点发散而只能逐个过程地重建；幂次修正的定义本身又与微扰级数的阶乘发散（重整化子）纠缠在一起。最后一节说明本书全部因子化定理作为关于强子的陈述处于什么地位。

\begin{tcolorbox}[enhanced,breakable,title=本章证明了什么／没有证明什么,fonttitle=\bfseries,
  colback=gray!3!white,colframe=gray!55!black,boxrule=0.6pt,arc=2pt,
  left=5pt,right=5pt,top=3pt,bottom=3pt,before skip=6pt,after skip=6pt]
\textbf{证明了：}\cref{prop:power-next-to-soft}（\tagR，$R_{\mathrm{PT}}$，树图层次）；\cref{prop:power-rearrangement}（\tagR，$R_{\mathrm{def}}$）；\cref{thm:power-ope-convergence}（\tagR，$R_{\mathrm{PT}}$，有质量欧氏 $\varphi^4_4$ 的每一阶；完整证明在文献中）；\cref{prop:power-adler-large-beta0}、\cref{prop:power-thooft-model}（\tagP：分别在大 $\beta_0$ 极限与二维大 $N_c$ 模型中完整证明）。\par
\textbf{给出全阶论证：}\cref{prop:power-dis-twist4-x}（\tagA，$x$ 空间；瓶颈：次领头幂次的区域分析只给出 E-S 型陈述，领头扭度部分继承\cref{thm:disall-factorization}的瓶颈，如 \cref{asm:appN-collins}）。\par
\textbf{没有证明：}一般的次领头幂次因子化（\cref{prop:power-nlp-endpoint}，\tagP）；强子—强子单自旋不对称的扭度 3 因子化（\cref{conj:power-twist3-pp}，\tagC）；红外重整化子与高扭度的一般对应（\cref{conj:power-renormalon-ht}，\tagC）；QCD 作为四维量子场论的构造与非微扰 OPE（\cref{rem:power-np-status}，\cref{asm:intro-wightman}、\cref{asm:intro-np-ope}）。
\end{tcolorbox}

\cref{sec:power-twist4,sec:power-twist3}讨论 DIS 的扭度 4 与单自旋不对称的扭度 3，\cref{sec:power-nlp}证明次软定理并说明端点发散，\cref{sec:power-ope-basis}证明 OPE 基的重排自由，\cref{sec:power-renormalons}在大 $\beta_0$ 极限下完整分析 Adler 函数的 Borel 平面，\cref{sec:power-np}给出非微扰地位与三个非微扰“实验室”。约定同全书（$D_\mu=\partial_\mu-\ii gA^a_\mu t^a$，$\MSbar$，$\mu^2\dd\as/\dd\mu^2=-\frac{\beta_0}{4\pi}\as^2+\dots$，$\beta_0=\frac{11}{3}C_A-\frac43T_Fn_f$）；$\Lambda$ 为 $\Lambda_{\MSbar}$（\cref{def:qcd-lambda}）或同量级的强子标度，$Q$ 为硬标度。

\begin{tcolorbox}[enhanced,breakable,title=常见误解（本章）,fonttitle=\bfseries,
  colback=yellow!3!white,colframe=orange!60!black,boxrule=0.6pt,arc=2pt,
  left=5pt,right=5pt,top=3pt,bottom=3pt,before skip=6pt,after skip=6pt]
\begin{enumerate}[label=(\arabic*),leftmargin=1.6em,itemsep=0pt,topsep=1pt]
\item % M10
\textbf{“幂次修正总是 $O(\Lambda^2/Q^2)$。”}事件形状的平均值有 $1/Q$ 修正\cite{DokshitzerWebber1995}，锥形喷注按 $1/(Q\delta)$\cite{Korchemsky:1994is}；扭度 3 的自旋不对称是 $O(M/Q)$；大 $x$ 处实际参数是 $\Lambda^2/[Q^2(1-x)]$（\cref{ex:power-twist4-largex}），核靶中有 $A^{1/3}$ 增强（\cref{prop:disall-nuclear}）。
\item % critic C11
\textbf{“余项是 $O(\Lambda^2/Q^2)$，因为高扭度算符的量纲高 2。”}这是 OPE 的陈述，对单举 DIS 与 $e^+e^-$ 的矩在非微扰 OPE 的前提下成立（\cref{asm:intro-np-ope}）；Drell--Yan 的扭度 4 广义因子化只有全阶论证（\cref{thm:dy-higher-twist}，\tagA），“没有 $\Lambda/Q$ 项”只有重整化子分析的证据（\cref{prop:power-renormalon-examples}）。
\item % M22
\textbf{“领头幂次因子化可以直接推广到次领头幂次。”}次领头幂次的卷积一般在端点发散，须逐个过程重新因子化（\cref{prop:power-nlp-endpoint}）；Drell--Yan 在 $m^4/s^2$ 处因子化已知失效（\cref{prop:dy-twist6-failure}，\tagV）。
\item % critic O20
\textbf{“重整化子与高扭度一一对应是已证明的。”}只在大 $\beta_0$ 极限与二维 $O(N)$ $\sigma$ 模型中确立，一般情形是\cref{conj:power-renormalon-ht}。
\item % M28 ; critic O21
\textbf{“[R] 是对真实 QCD 的证明；强子层次的因子化有非微扰证明。”}QCD 尚未被构造，[R] 总带限定语（\cref{def:intro-tags}）；非微扰“实验室”只有微扰 $\varphi^4_4$ 中 OPE 的收敛、't Hooft 模型与格点证据（\cref{sec:power-np}）。
\end{enumerate}
\end{tcolorbox}
